Line $MN$ is $2x+3y=6$, whereas $CB$ is $4x+3y=24$, so $K=(9,-4)$. The two circle equations are
\[
x^2+y^2-3x-2y=0,\qquad x^2+y^2-6x-8y=0.
\]
Their common chord lies on $x+2y=0$. Substitution gives $S=\left(\dfrac{8}{5},-\dfrac{4}{5}\right)$.

For line $SM$, write $(x,y)=S+t(7,4)$. Substitution into the second circle equation reduces to $t(65t-58)=0$, so
\[
P=\left(\dfrac{102}{13},\dfrac{36}{13}\right).
\]
For line $SK$, write $(x,y)=S+t(37,-16)$. Substitution gives $t(1625t+50)=0$, so
\[
Q=\left(\dfrac{6}{13},-\dfrac{4}{13}\right).
\]
Consequently
\[
PQ=\sqrt{\left(\dfrac{96}{13}\right)^2+\left(\dfrac{40}{13}\right)^2}=8=AC.
\]
\begin{asy}
texpreamble("\usepackage{amsmath,amssymb}");
size(240);
defaultpen(fontsize(9pt));
pair A=(0,0),B=(6,0),C=(0,8),M=(3,0),N=(0,2),K=(9,-4),S=(8/5,-4/5),P=(102/13,36/13),Q=(6/13,-4/13);draw(circle((3,4),5));draw(circle((1.5,1),sqrt(13)/2),lightblue);draw(A--B--C--cycle);draw(N--K,gray);draw(S--P,blue);draw(Q--K,blue);draw(P--Q,red);label("$A$",A,W);label("$B$",B,NE);label("$C$",C,N);label("$M$",M,N);label("$N$",N,W);label("$K$",K,SE);label("$S$",S,S);label("$P$",P,E);label("$Q$",Q,SW);
\end{asy}
